\documentclass[11pt]{article}
\usepackage[margin=1in]{geometry}
\usepackage{amsmath,amssymb,mathtools}
\usepackage{fontspec,unicode-math}
\setmainfont{Latin Modern Roman}
\setsansfont{Latin Modern Sans}
\setmathfont{Latin Modern Math}
\usepackage{xurl,hyperref}
\hypersetup{colorlinks=true,linkcolor=blue,urlcolor=blue}
\newcommand{\sref}[2]{\hyperlink{source:#1:#2}{[S#2]}}
\newcommand{\cataloguescope}[1]{\paragraph{Mathematical question.}#1}
\begin{document}
\section{Local connectivity of geometrically finite Julia components}\label{statement}

\subsection{Definitions and mathematical statement}

\paragraph{Shared notation.}$\mathbb N=\{1,2,\ldots\}$, and $\mathbb Z,\mathbb Q,\mathbb R,\mathbb C$ denote the integers, rationals, reals and complex numbers. Any other conventions are specified locally.


For a rational map $f:\widehat{\mathbb C}\to\widehat{\mathbb C}$ of degree at least two, $J(f)$ is the complement of its Fatou set, where the iterates form a normal family. If $C(f)$ is its set of critical points, set $P(f)=\bigcup_{n\geq1}f^n(C(f))$. In the source convention $f$ is geometrically finite when $\overline{P(f)}\cap J(f)$ is finite; parabolic cycles are allowed.A space is locally connected if every point has a neighborhood basis of connected open subsets. Components here are connected components with the subspace topology from the sphere.
\paragraph{Statement recorded in the supplied source.}
Prove that every connected component of $J(f)$ is locally connected. \sref{statement}{2}
\subsection{Short English statement}
Prove local connectivity of every Julia component of a geometrically finite rational map.
\subsection{Sources}
\begin{description}
\item[\hypertarget{source:statement:1}{S1}] ULAM.AI and the UnsolvedMath Contributors, UnsolvedMath: A Curated Collection of Open Mathematics Problems, catalogue entry. CC BY 4.0. \url{https://huggingface.co/datasets/ulamai/UnsolvedMath}.
\item[\hypertarget{source:statement:2}{S2}] Ben Bielefeld (editor) and conference contributors, \emph{Conformal Dynamics Problem List}, arXiv:math/9201271v1 (1990), Curt McMullen, \S5, Geometrically finite rational maps, Problem 2, printed p. 11. \url{https://arxiv.org/abs/math/9201271v1}.
\end{description}
\label{end:statement}
\end{document}
